\documentclass[12pt, a4paper]{article}

% ---------- packages ----------
\usepackage[T1]{fontenc}
\usepackage[utf8]{inputenc}

\usepackage[margin=2.5cm]{geometry}
\usepackage{amsmath, amsthm, amssymb}
\usepackage{mathtools}
\usepackage{xcolor}
\usepackage{booktabs}
\usepackage{array}
\usepackage{hyperref}
\usepackage{tcolorbox}
\usepackage{mdframed}
\usepackage{parskip}
\usepackage{enumitem}
\usepackage{titlesec}

% ---------- colour palette ----------
\definecolor{accentblue}{RGB}{30, 60, 110}
\definecolor{lightblue}{RGB}{220, 232, 248}
\definecolor{warmgray}{RGB}{245, 242, 237}
\definecolor{borderblue}{RGB}{80, 120, 180}
\definecolor{notegreen}{RGB}{0, 100, 60}
\definecolor{darkorange}{RGB}{170, 80, 0}

% ---------- section formatting ----------
\titleformat{\section}{\large\bfseries\color{accentblue}}{}{0em}{\thesection\quad}[\vspace{-0.3em}\textcolor{accentblue}{\rule{\linewidth}{0.8pt}}]
\titleformat{\subsection}{\normalsize\bfseries\color{accentblue}}{}{0em}{\thesubsection\quad}

% ---------- tcolorbox styles ----------
\tcbuselibrary{breakable, skins}

\newtcolorbox{keybox}[1][]{
  enhanced, breakable,
  colback=lightblue, colframe=accentblue,
  boxrule=1pt, arc=4pt,
  fonttitle=\bfseries\small\color{accentblue},
  title={#1}
}

\newtcolorbox{resultbox}[1][]{
  enhanced, breakable,
  colback=warmgray, colframe=accentblue,
  boxrule=1.5pt, arc=4pt,
  fonttitle=\bfseries\small\color{accentblue},
  title={#1}
}

\newtcolorbox{notebox}[1][]{
  enhanced, breakable,
  colback=white, colframe=darkorange,
  boxrule=0.8pt, leftrule=4pt, arc=2pt,
  fonttitle=\bfseries\small\color{darkorange},
  title={#1}
}

% ---------- custom commands ----------
\newcommand{\E}{\mathbb{E}}
\newcommand{\Var}{\operatorname{Var}}
\newcommand{\dd}{\,\mathrm{d}}
\newcommand{\kap}{\kappa}

% ---------- title ----------
\title{%
  \textbf{\textcolor{accentblue}{Liquidating Two Correlated Assets}}\\[4pt]
  \large Almgren--Chriss Framework with Cross-Asset Hedging\\
  \normalsize Mean--Variance Criterion, Bellman Backward Induction\\[6pt]
}
\author{}
\date{}

% ============================================================
\begin{document}
\maketitle
\thispagestyle{empty}

% -------------------------------------------------------------------
\section{Model Setup}
% -------------------------------------------------------------------

We now extend the single-asset Almgren--Chriss framework to two risky assets
$S^1$ and $S^2$ with correlated price dynamics.

\paragraph{Price dynamics.}
Both assets follow arithmetic Brownian motions (no drift, no permanent impact):
\[
  \dd S^1_t = \sigma_1 \dd W^1_t,
  \qquad
  \dd S^2_t = \sigma_2 \dd W^2_t,
  \qquad
  \dd\langle W^1, W^2\rangle_t = \rho\,\dd t,
\]
where $\sigma_1, \sigma_2 > 0$ are volatilities and $\rho \in (-1,1)$ is the
instantaneous correlation between the two Brownian motions.

\paragraph{Inventory and constraints.}
\begin{itemize}[leftmargin=1.5em]
  \item $q^1_t$: shares of $S^1$ held at time $t$, with $q^1_0 = Q > 0$ and $q^1_T = 0$.
  \item $q^2_t$: shares of $S^2$ held at time $t$, with $q^2_0 = 0$ and $q^2_T = 0$.
  \item $v^i_t = -\dot{q}^i_t$: trading speed for asset $i \in \{1,2\}$.
\end{itemize}
The trader starts with $Q$ shares of $S^1$ and no position in $S^2$.
Both positions must be flat at time $T$.

\paragraph{Temporary market impact.}
Trading at speed $v^i_t$ incurs a quadratic liquidity cost:
\[
  \text{impact cost} = \eta_i (v^i_t)^2 \dd t,
\]
where $\eta_1, \eta_2 > 0$ are the impact coefficients.
We assume $\eta_2 < \eta_1$: asset $S^2$ is more liquid than $S^1$.

\begin{notebox}[Why trade $S^2$ at all?]
  Although the trader holds no $S^2$ initially, trading it is useful when
  $|\rho|$ is large: going short $S^2$ (if $\rho > 0$) or long $S^2$
  (if $\rho < 0$) creates a hedge that offsets the price risk of holding $S^1$,
  at a lower liquidity cost since $\eta_2 \ll \eta_1$.
\end{notebox}

% -------------------------------------------------------------------
\section{Cash Flow and Payoff}
% -------------------------------------------------------------------

\paragraph{Cash account dynamics.}
The trader's cash account $X_t$ evolves as:
\[
  \dd X_t
  = q^1_t \dd S^1_t + q^2_t \dd S^2_t
  - \eta_1 (v^1_t)^2 \dd t - \eta_2 (v^2_t)^2 \dd t.
\]
The first two terms are the mark-to-market P\&L from holding the positions;
the last two are the liquidity costs from trading.

\paragraph{Integrated cash flow.}
Integrating over $[0,T]$:
\[
  X_T = X_0
  + \int_0^T q^1_t \dd S^1_t
  + \int_0^T q^2_t \dd S^2_t
  - \eta_1\int_0^T (v^1_t)^2 \dd t
  - \eta_2\int_0^T (v^2_t)^2 \dd t.
\]

\paragraph{IS payoff.}
We measure performance relative to the IS benchmark $QS^1_0$
(the ideal cash received if all shares of $S^1$ were sold instantly at $t=0$):
\[
  \Pi = X_T - QS^1_0.
\]
By the Almgren--Chriss normalisation $X_0 = QS^1_0$, this simplifies to:
\[
  \Pi
  = \int_0^T q^1_t \dd S^1_t
  + \int_0^T q^2_t \dd S^2_t
  - \eta_1\int_0^T (v^1_t)^2 \dd t
  - \eta_2\int_0^T (v^2_t)^2 \dd t.
\]

% -------------------------------------------------------------------
\section{Mean--Variance Criterion}
% -------------------------------------------------------------------

\subsection{Objective}

Following the course formalism, we maximise:
\[
  \max_{v^1,\,v^2}\;\E[\Pi] - \lambda\,\Var(\Pi),
\]
or equivalently minimise:
\[
  \min_{v^1,\,v^2}\;-\E[\Pi] + \lambda\,\Var(\Pi).
\]

\subsection{Deriving \texorpdfstring{$-\E[\Pi]$}{-E[Pi]}}

Since strategies are deterministic and $\dd S^i_t = \sigma_i \dd W^i_t$ is a
martingale, the stochastic integrals have zero expectation:
\[
  \E\!\left[\int_0^T q^i_t \dd S^i_t\right] = 0, \qquad i = 1, 2.
\]
Therefore:
\[
  \E[\Pi] = -\eta_1\int_0^T (v^1_t)^2 \dd t - \eta_2\int_0^T (v^2_t)^2 \dd t,
\]
and:
\begin{equation}
  -\E[\Pi] = \eta_1\int_0^T (v^1_t)^2 \dd t + \eta_2\int_0^T (v^2_t)^2 \dd t.
  \label{eq:neg-mean}
\end{equation}

\subsection{Deriving \texorpdfstring{$\Var(\Pi)$}{Var(Pi)}}

The payoff decomposes into a deterministic part (the liquidity costs) and a
stochastic part (the P\&L from price moves). Since the deterministic part
does not contribute to variance:
\[
  \Var(\Pi)
  = \E\!\left[
      \left(\int_0^T q^1_t \dd S^1_t + \int_0^T q^2_t \dd S^2_t\right)^{\!2}
    \right].
\]
Expanding the square:
\[
  \Var(\Pi)
  = \E\!\left[\left(\int_0^T q^1_t \dd S^1_t\right)^{\!2}\right]
  + \E\!\left[\left(\int_0^T q^2_t \dd S^2_t\right)^{\!2}\right]
  + 2\,\E\!\left[\int_0^T q^1_t \dd S^1_t \cdot \int_0^T q^2_t \dd S^2_t\right].
\]

\paragraph{Diagonal terms (It\^{o} isometry).}
Since $\dd S^i_t = \sigma_i \dd W^i_t$ and $q^i_t$ is deterministic:
\[
  \E\!\left[\left(\int_0^T q^i_t \dd S^i_t\right)^{\!2}\right]
  = \sigma_i^2 \int_0^T (q^i_t)^2 \dd t.
\]

\paragraph{Cross term.}
Using the quadratic covariation $\dd\langle S^1, S^2\rangle_t = \rho\sigma_1\sigma_2\,\dd t$:
\[
  \E\!\left[\int_0^T q^1_t \dd S^1_t \cdot \int_0^T q^2_t \dd S^2_t\right]
  = \rho\sigma_1\sigma_2\int_0^T q^1_t q^2_t \dd t.
\]

\paragraph{Full variance.}
\begin{equation}
  \Var(\Pi)
  = \sigma_1^2\int_0^T (q^1_t)^2 \dd t
  + \sigma_2^2\int_0^T (q^2_t)^2 \dd t
  + 2\rho\sigma_1\sigma_2\int_0^T q^1_t q^2_t \dd t.
  \label{eq:variance}
\end{equation}

\begin{keybox}[Role of the cross term]
  The cross term $2\rho\sigma_1\sigma_2\int_0^T q^1_t q^2_t\,\dd t$ is the key
  to understanding the hedge:
  \begin{itemize}[leftmargin=1.5em]
    \item If $\rho > 0$ and the trader goes \textbf{short} $S^2$ ($q^2_t < 0$),
          then $q^1_t q^2_t < 0$, making the cross term negative — reducing variance.
    \item If $\rho < 0$ and the trader goes \textbf{long} $S^2$ ($q^2_t > 0$),
          then $q^1_t q^2_t > 0$ but $\rho < 0$, again making the cross term negative.
    \item In both cases, the sign of the optimal $q^2$ follows $-\operatorname{sign}(\rho)$:
          the hedge direction is always set to reduce variance.
  \end{itemize}
\end{keybox}

\subsection{Full Optimisation Problem}

Combining \eqref{eq:neg-mean} and \eqref{eq:variance}:

\begin{resultbox}[Two-asset mean--variance problem]
\[
  \min_{v^1,\,v^2}\;
  \eta_1\int_0^T (v^1_t)^2 \dd t
  + \eta_2\int_0^T (v^2_t)^2 \dd t
  + \lambda\!\left(
      \sigma_1^2\int_0^T (q^1_t)^2 \dd t
      + \sigma_2^2\int_0^T (q^2_t)^2 \dd t
      + 2\rho\sigma_1\sigma_2\int_0^T q^1_t q^2_t \dd t
    \right)
\]
subject to $q^1_0 = Q,\; q^1_T = 0,\; q^2_0 = 0,\; q^2_T = 0$.
\end{resultbox}

\begin{notebox}[Generalisation to $n$ assets]
  With $n$ assets, covariance matrix $\Sigma_{ij} = \rho_{ij}\sigma_i\sigma_j$,
  and diagonal impact matrix $H = \mathrm{diag}(\eta_1,\ldots,\eta_n)$, the
  problem takes the compact matrix form:
  \[
    \min_{\mathbf{v}}
    \int_0^T \mathbf{v}_t^\top H\,\mathbf{v}_t \dd t
    + \lambda\int_0^T \mathbf{q}_t^\top \Sigma\,\mathbf{q}_t \dd t.
  \]
  The 2-asset case above is the special case $n = 2$.
\end{notebox}

% -------------------------------------------------------------------
\section{Why No Closed-Form Solution?}
% -------------------------------------------------------------------

In the single-asset case, the Euler--Lagrange equation yields a single ODE
$\ddot{q} = \kap^2 q$ with explicit $\sinh/\cosh$ solution.

For two assets, the optimality conditions give a \emph{coupled} system:
\[
  \begin{cases}
    \eta_1\,\ddot{q}^1 = \lambda\bigl(\sigma_1^2\,q^1 + \rho\sigma_1\sigma_2\,q^2\bigr), \\[4pt]
    \eta_2\,\ddot{q}^2 = \lambda\bigl(\sigma_2^2\,q^2 + \rho\sigma_1\sigma_2\,q^1\bigr).
  \end{cases}
\]
This is a \textbf{Boundary Value Problem (BVP)}: conditions are imposed at
\emph{both} ends of the interval,
\[
  q^1(0) = Q,\; q^2(0) = 0 \qquad \text{and} \qquad q^1(T) = 0,\; q^2(T) = 0.
\]
To integrate forward from $t=0$, one would also need $\dot{q}^1(0)$ and
$\dot{q}^2(0)$, which are \emph{unknown} and depend on all parameters
$(\rho, \eta_1, \eta_2, \sigma_1, \sigma_2)$ in a non-trivial way.
For general parameters, this system admits no clean closed-form solution.

\begin{keybox}[When does a closed form exist?]
  In the special case $\sigma_1 = \sigma_2$ and $\eta_1 = \eta_2$, the two
  equations decouple after a rotation of coordinates, and explicit
  $\sinh/\cosh$ solutions can be recovered. In the general case,
  Bergault, Drissi and Gu\'eant (2021) show that the solution satisfies a
  \emph{matrix Riccati ODE}, which requires numerical integration.
\end{keybox}

% -------------------------------------------------------------------
\section{Bellman Backward Induction on a 2D Grid}
\label{sec:bellman}
% -------------------------------------------------------------------

\subsection{Why Bellman Works Here}

Bellman backward induction sidesteps the BVP entirely:
\begin{itemize}[leftmargin=1.5em]
  \item Work \textbf{backwards from $T$}: the terminal condition
        $V_N(0,0) = 0$ and $V_N(q^1,q^2) = +\infty$ for
        $(q^1,q^2) \neq (0,0)$ is trivially known.
  \item At each earlier step $k$, minimise over the \emph{finite} set of
        feasible next states — no need to guess initial velocities.
  \item The terminal constraint $q^1(T) = q^2(T) = 0$ is enforced
        exactly by the boundary condition.
\end{itemize}

\subsection{Discrete Bellman Equation}

Let $\Delta t = T/N$. At each time step $k \in \{0,\ldots,N-1\}$, the trader
chooses $(n_1, n_2)$ — the number of shares of $S^1$ sold and the change in
$S^2$ position — to minimise the sum of immediate cost and future cost-to-go:

\begin{resultbox}[Two-asset Bellman equation]
\[
  V_k(q^1, q^2)
  = \min_{\substack{0 \le n_1 \le q^1 \\ n_2 \in \mathbb{R}}}
  \left[
    \eta_1\!\left(\frac{n_1}{\Delta t}\right)^{\!2}\!\Delta t
    + \eta_2\!\left(\frac{n_2}{\Delta t}\right)^{\!2}\!\Delta t
    + \lambda\Bigl(\sigma_1^2 (q^1)^2 + \sigma_2^2 (q^2)^2
                   + 2\rho\sigma_1\sigma_2\,q^1 q^2\Bigr)\Delta t
    + V_{k+1}(q^1 - n_1,\; q^2 - n_2)
  \right]
\]
with $V_N(0,0) = 0$ and $V_N(q^1,q^2) = +\infty$ for $(q^1,q^2)\neq(0,0)$.
\end{resultbox}

\paragraph{Term-by-term reading.}
\begin{itemize}[leftmargin=1.5em]
  \item $\eta_1(n_1/\Delta t)^2\Delta t$: liquidity cost of selling $n_1$ shares
        of $S^1$ at velocity $n_1/\Delta t$.
  \item $\eta_2(n_2/\Delta t)^2\Delta t$: liquidity cost of trading $n_2$ shares
        of $S^2$.
  \item $\lambda(\sigma_1^2(q^1)^2 + \sigma_2^2(q^2)^2 +
        2\rho\sigma_1\sigma_2\,q^1 q^2)\Delta t$: instantaneous inventory risk,
        including the cross term that enables hedging.
  \item $V_{k+1}(q^1-n_1, q^2-n_2)$: future optimal cost-to-go.
\end{itemize}

\paragraph{Constraints on $n_2$.}
The action $n_2$ is unrestricted in sign:
\begin{itemize}[leftmargin=1.5em]
  \item $n_2 > 0$: increasing a short position in $S^2$ (or reducing a long).
  \item $n_2 < 0$: covering a short / going long $S^2$.
\end{itemize}
The optimal sign of $q^2_t$ is determined endogenously by the optimisation:
\begin{itemize}[leftmargin=1.5em]
  \item $\rho > 0 \Rightarrow$ optimal $q^2_t < 0$ (short hedge).
  \item $\rho < 0 \Rightarrow$ optimal $q^2_t > 0$ (long hedge).
  \item $\rho = 0 \Rightarrow$ optimal $q^2_t = 0$ (no hedge benefit).
\end{itemize}

\paragraph{Computational complexity.}
The state space is $O(n_q^2)$ and actions are $O(n_q^2)$ per state,
giving total cost $O(N \cdot n_q^4)$. Keep $n_q$ small (10--30) for
tractability; a finer grid reduces discretisation artefacts (staircasing)
at the cost of computation time.

\begin{notebox}[Sanity check]
  Setting $q^2 = 0$ and $\sigma_1=\sigma$, $\eta_1=\eta$ recovers exactly
  the single-asset IS Bellman equation: the cross term vanishes, the $\eta_2$
  term drops out, and $V_k(q^1) = \min_{0\le n\le q^1}[\eta(n/\Delta t)^2\Delta t
  + \lambda\sigma^2(q^1)^2\Delta t + V_{k+1}(q^1-n)]$.
\end{notebox}

% -------------------------------------------------------------------
\section{Question 2: Impact of the Hedge on the Optimal Strategy}
\label{sec:impact}
% -------------------------------------------------------------------

We compare two strategies:
\begin{itemize}[leftmargin=1.5em]
  \item \textbf{Single-asset IS}: the trader can only sell $S^1$,
        using the optimal strategy from Section 1.
  \item \textbf{Two-asset}: the trader can also trade $S^2$,
        using the 2D Bellman above.
\end{itemize}

Both strategies liquidate $S^1$ from $Q$ to $0$ over $[0,T]$.
The comparison focuses on two quantities:

\paragraph{Portfolio variance.}
Using formula \eqref{eq:variance} (with $q^2 \equiv 0$ for single-asset IS):
\[
  \Var_{\mathrm{single}} = \sigma_1^2\int_0^T (q^1_t)^2 \dd t,
\]
\[
  \Var_{\mathrm{two}} = \sigma_1^2\int_0^T (q^1_t)^2 \dd t
  + \sigma_2^2\int_0^T (q^2_t)^2 \dd t
  + 2\rho\sigma_1\sigma_2\int_0^T q^1_t q^2_t \dd t.
\]

\paragraph{Liquidity cost.}
\[
  \mathrm{Cost}_{\mathrm{single}} = \eta_1\int_0^T (v^1_t)^2 \dd t,
  \qquad
  \mathrm{Cost}_{\mathrm{two}} = \eta_1\int_0^T (v^1_t)^2 \dd t
  + \eta_2\int_0^T (v^2_t)^2 \dd t.
\]

\begin{resultbox}[Key findings]
\begin{enumerate}[leftmargin=1.5em]
  \item \textbf{Variance reduction.}
        The cross term $2\rho\sigma_1\sigma_2\int q^1 q^2\,\dd t$ is negative
        in the optimal two-asset strategy (since $q^1 > 0$ and $\rho q^2 < 0$),
        directly reducing $\Var_{\mathrm{two}}$ below $\Var_{\mathrm{single}}$.
        With $\rho = 0.95$ and $\eta_2/\eta_1 = 0.05$, the numerical results
        show a variance reduction of approximately $51\%$.

  \item \textbf{Modest additional cost.}
        Since $\eta_2 \ll \eta_1$, shorting $S^2$ is cheap.
        $\mathrm{Cost}_{\mathrm{two}} \approx \mathrm{Cost}_{\mathrm{single}}$
        in practice — the hedge pays for itself.

  \item \textbf{Slower, quasi-linear liquidation of $S^1$.}
        With the hedge absorbing inventory risk, the urgency to front-load
        $S^1$ sales disappears. The optimal $q^1_t$ is quasi-linear
        (close to TWAP) rather than the convex IS curve. This reflects the
        \emph{decoupling} of risk management (delegated to $S^2$) from
        liquidation (spread evenly over $[0,T]$).

  \item \textbf{Symmetry in $\rho$.}
        The speed of $S^1$ liquidation depends only on $|\rho|$, not its sign.
        What changes with the sign of $\rho$ is only the \emph{direction}
        of the $S^2$ position: short for $\rho > 0$, long for $\rho < 0$.
\end{enumerate}
\end{resultbox}

\subsection{Economic Intuition for the $S^2$ Path}

The optimal $q^2_t$ follows a characteristic arc:
\begin{enumerate}[leftmargin=1.5em]
  \item \textbf{Early phase}: $q^1_t$ is large, so inventory risk is high.
        The trader shorts $S^2$ aggressively to offset $S^1$ exposure via
        the cross term. This is cheap since $\eta_2 \ll \eta_1$.
  \item \textbf{Middle phase}: the marginal benefit of further hedging
        (proportional to $q^1_t$) equals the marginal cost
        (proportional to $\eta_2 v^2_t$). The short position stabilises.
  \item \textbf{Late phase}: as $q^1_t \to 0$, the cross term $\propto q^1_t q^2_t$
        vanishes, but holding a short in $S^2$ still costs $\sigma_2^2(q^2_t)^2\lambda$.
        The trader covers the short to meet the constraint $q^2_T = 0$.
\end{enumerate}

The minimum of $q^2_t$ (deepest short) occurs when
$\lambda\rho\sigma_1\sigma_2\,q^1_t = \eta_2\,v^2_t$:
the marginal hedge benefit exactly equals the marginal liquidity cost.

% -------------------------------------------------------------------
\section{Summary}
% -------------------------------------------------------------------

\begin{resultbox}[Two-asset liquidation: key results]
\begin{center}
  \renewcommand{\arraystretch}{1.6}
  \begin{tabular}{@{}lll@{}}
    \toprule
    \textbf{Feature} & \textbf{Single-asset IS} & \textbf{Two-asset} \\
    \midrule
    $q^1_t$ shape & Convex (front-loaded) & Quasi-linear \\
    $q^2_t$ & $0$ & Short ($\rho>0$) or Long ($\rho<0$) \\
    Variance & $\sigma_1^2\int(q^1)^2\,\dd t$ & Reduced by cross term \\
    Cost & $\eta_1\int(v^1)^2\,\dd t$ & $+\,\eta_2\int(v^2)^2\,\dd t$ (small) \\
    Solution method & Closed form & Bellman on 2D grid \\
    \bottomrule
  \end{tabular}
\end{center}

\medskip
The hedge is most effective when $|\rho|$ is large (assets move together)
and $\eta_2 \ll \eta_1$ (trading $S^2$ is cheap).
The sign of the optimal $q^2$ always satisfies
$\operatorname{sign}(q^2_t) = -\operatorname{sign}(\rho)$.
\end{resultbox}

\end{document}
